\documentclass[11pt,a4paper]{article}

\usepackage{kotex}
\usepackage{fontspec}
\setmainfont{Times New Roman}
\setmainhangulfont{AppleMyungjo}
\setsanshangulfont{Apple SD Gothic Neo}

\usepackage[a4paper,top=18mm,bottom=18mm,left=20mm,right=20mm]{geometry}
\usepackage{fancyhdr}
\setlength{\headheight}{24pt}
\usepackage{enumitem}
\usepackage{mdframed}
\usepackage{array}
\usepackage{booktabs}

\setlength{\parindent}{1em}
\linespread{1.18}

\pagestyle{fancy}
\fancyhf{}
\renewcommand{\headrulewidth}{0.6pt}
\fancyhead[L]{\sffamily\bfseries 2026 고1 영어 변형 --- 정답 및 해설 21}
\fancyhead[R]{\sffamily [shop for words]}
\renewcommand{\footrulewidth}{0pt}
\fancyfoot[C]{\framebox[16mm][c]{\thepage}}

\newcommand{\qno}[1]{\par\vspace{10pt}\noindent{\sffamily\bfseries #1.}\ }
\newmdenv[linewidth=0.6pt,roundcorner=3pt,innertopmargin=7pt,innerbottommargin=7pt,
          innerleftmargin=9pt,innerrightmargin=9pt,skipabove=5pt,skipbelow=5pt]{pbox}

\begin{document}

\begin{center}
{\fontsize{15}{17}\selectfont\sffamily\bfseries [지문 21] 정답 및 해설}
\end{center}
\vspace{6pt}

%% 정답 표
\begin{center}
\begin{tabular}{cccc}
\toprule
\textbf{문항} & \textbf{유형} & \textbf{정답} & \textbf{배점} \\
\midrule
1 & 빈칸추론 & ③ & 3점 \\
2 & 어법 (네모) & ② & 2점 \\
3 & 서술형 --- 해석 & 하단 모범답안 참조 & 3점 \\
4 & 영어 제목 & ② & 2점 \\
\bottomrule
\end{tabular}
\end{center}

\bigskip

%% 문항 1 해설
\qno{1}\ \textbf{[빈칸추론 --- 3점] 정답: ③ rationalizations}

\vspace{4pt}
\noindent\textbf{근거 문장}: \textit{``But in practice, most of the time we are in a market for \underline{\hspace{3.5cm}}. We shop for words not because they contain ideas, but because they contain stories about ideas. \ldots\ What we seek are statements to justify those beliefs --- the lawyer's way.''}

\vspace{4pt}
\noindent\textbf{해설}: 글쓴이는 이상적으로는 사실에 가장 잘 맞는 아이디어를 평가하는 ``과학자의 방식''을 언급하지만, 현실에서는 우리가 이미 믿고 있는 것을 정당화해 줄 말을 찾는다고 대조한다. 빈칸 뒤에 \textit{``the lawyer's way''}라는 단서가 이 사후 합리화(post-hoc rationalization) 개념을 직접 가리킨다. 따라서 \textbf{rationalizations}(사후 합리화, 정당화)가 가장 적절하다.

\vspace{6pt}
\noindent\textbf{오답 소거}:
\begin{itemize}[leftmargin=2em,itemsep=1pt,topsep=2pt]
  \item ① inspirations --- 영감은 새로운 아이디어를 찾는 것이므로 지문의 ``already know what we believe''와 모순.
  \item ② innovations --- 혁신을 구매하는 시장이라는 개념은 지문 논지와 무관.
  \item ④ contradictions --- 모순을 찾는다는 의미는 지문 흐름(기존 믿음의 강화)에 반함.
  \item ⑤ observations --- 관찰·사실을 구한다면 ``과학자의 방식''에 해당하므로 현실 묘사와 불일치.
\end{itemize}

\bigskip

%% 문항 2 해설
\qno{2}\ \textbf{[어법 네모 --- 2점] 정답: ②\quad (A) fit \quad (B) form \quad (C) better}

\vspace{4pt}
\noindent\textbf{(A) fit}: \textit{``their best fit to reality''} --- \textbf{fit}은 명사로 ``적합(성)''을 뜻하며, \textit{best}의 수식을 받는 명사 자리이다. \textit{fitting}은 형용사·분사이므로 \textit{their best fitting to reality}는 비문.

\vspace{4pt}
\noindent\textbf{(B) form}: \textit{``let the best and brightest facts form the basis''} --- 사역동사 \textbf{let}은 목적격 보어로 \textbf{동사원형(bare infinitive)}을 취한다. \textit{to form}은 오답.

\vspace{4pt}
\noindent\textbf{(C) better}: \textit{``we can use it to do better''} --- \textit{do better}은 숙어적 표현으로 ``더 잘하다''를 의미하며, 앞 절(\textit{good for this function})과 대조하는 비교급이 필요하다. \textit{good}은 형용사이므로 \textit{do good}(선행을 하다)의 다른 의미가 되어 문맥 불일치.

\bigskip

%% 문항 3 해설
\qno{3}\ \textbf{[서술형 --- 해석 --- 3점]}

\vspace{4pt}
\noindent\textbf{대상 문장}: \textit{By treating language with greater respect, we can achieve more than simple self-defense.}

\vspace{6pt}
\begin{pbox}
\textbf{모범답안}: 언어를 더욱 존중하여 다룸으로써, 우리는 단순한 자기 방어 이상을 달성할 수 있다.
\end{pbox}

\vspace{6pt}
\noindent\textbf{채점 포인트}:
\begin{itemize}[leftmargin=2em,itemsep=1pt,topsep=2pt]
  \item \textit{By treating language with greater respect} --- ``언어를 더욱(더 크게) 존중하여 다룸으로써'' (1점)
  \item \textit{we can achieve more than} --- ``우리는 \ldots\ 이상을 달성할 수 있다'' (1점)
  \item \textit{simple self-defense} --- ``단순한 자기 방어'' (1점)
\end{itemize}

\vspace{4pt}
\noindent\textit{* ``By treating $\sim$''의 동명사 의미상 주어(we)와 주절 주어 일치 확인. ``more than''을 ``보다 많이''로 직역해도 문맥상 수용 가능.}

\vspace{8pt}
\noindent{\sffamily\bfseries 4. 제목 - 정답: ②}
\par 글은 사람들이 진리를 찾기보다 이미 가진 믿음을 정당화할 말을 찾는 경향을 비판한다. 따라서 언어가 기존 믿음을 방어하는 도구로 쓰인다는 ②가 글의 핵심을 가장 잘 담는다.

\end{document}
